\chapter{Introduction}
\label{ch:introduction}

\section{Motivation}
\label{sec:motivation}

The contemporary digital ecosystem is characterized by an accelerating proliferation
of distributed systems that have become the de facto standard for modern software
delivery \cite{coulouris2011distributed}. Primarily instantiated through
microservice-based deployments \cite{newman2015building, dragoni2017microservices},
this paradigm decomposes monolithic applications into independently deployable,
loosely coupled services. While this delivers substantial benefits in development
velocity, scalability, and fault isolation, it simultaneously introduces
unprecedented complexity in system observability and anomaly detection
\cite{kalloniatis2020microservices_observability}. Traditional centralized monitoring
paradigms, adequate for monolithic architectures, are rendered fundamentally
insufficient in this context \cite{lin2018aiops}.

\section{Problem Statement}
\label{sec:problem_statement}

Detecting anomalies across highly decentralized deployments introduces significant
technical challenges. The sheer volume of telemetry generated by thousands of
concurrent service instances creates a data ingestion bottleneck
\cite{sigelman2010dapper}, while the ephemeral nature of cloud-native environments
renders static monitoring configurations obsolete \cite{lin2018aiops}. Maintaining
strict service level objectives (SLOs) in this volatile environment necessitates a
paradigm shift: from reactive, centralized aggregation toward proactive,
privacy-preserving, and intelligent distributed monitoring \cite{lin2018aiops}.

\section{Research Objectives}
\label{sec:research_objectives}

The primary objective of this research is to design and evaluate a decentralized
architecture for anomaly detection in distributed systems that guarantees operational
visibility without compromising data privacy. This breaks down into five concrete
sub-objectives:

\begin{enumerate}
    \item \textbf{On-Device Model Architecture}: Design a lightweight LSTM-Autoencoder
    (LSTM-AE) \cite{malhotra2015lstm} for localized anomaly detection at the edge.

    \item \textbf{Privacy-Preserving Collaborative Training}: Integrate federated
    learning \cite{mcmahan2017communication} with differential privacy guarantees
    \cite{dwork2006differential, abadi2016deep} for robust global model training
    without exposing local telemetry.

    \item \textbf{Intelligent Root Cause Identification}: Construct a root cause
    analysis engine leveraging causal-dependency graphs to contextualize metrics and
    identify anomalous service interactions.

    \item \textbf{Contextual Anomaly Thresholding}: Develop a Q-learning-driven
    \cite{watkins1992qlearning} adaptive thresholding mechanism that optimizes
    anomaly sensitivity on a per-service basis.

    \item \textbf{Optimized Infrastructure Communication}: Implement and benchmark
    model weight serialization pipelines contrasting Protocol Buffers with Zstandard
    and LZ4 compression \cite{collet2018zstandard, lz4_reference}.
\end{enumerate}

\section{Research Contributions}
\label{sec:contributions}

This thesis delivers the following novel contributions to distributed systems
observability and AIOps \cite{lin2018aiops}:

\begin{itemize}
    \item \textbf{Decentralized Anomaly Framework}: A federated LSTM-Autoencoder
    framework tailored for heterogeneous distributed environments.

    \item \textbf{Secure Aggregation Pipeline}: A differential-privacy-aware FedAvg
    \cite{mcmahan2017communication} aggregation pipeline over operational telemetry.

    \item \textbf{Topological Root Cause Scoring}: A causal-dependency graph approach
    with a PageRank-inspired heuristic \cite{page1999pagerank} for scalable fault
    attribution.

    \item \textbf{Reinforcement Learning for Thresholds}: Q-learning applied to
    autonomously manage adaptive thresholding, minimizing false positives and
    negatives per service.

    \item \textbf{Communication Cost Empirics}: Empirical evaluation of LZ4 and Zstd
    compression schemes in bandwidth-constrained federated deployments.
\end{itemize}
